% Ready-to-input section; bibliography entries are supplied in the companion note.
\subsection{Real weights over global fields}
For a number field $K$, use one place for each real embedding and each
conjugate pair of complex embeddings, and set
\[
 L_v(a)=
 \begin{cases}
  \log|\sigma_v(a)|,&v\text{ real},\\
  2\log|\sigma_v(a)|,&v\text{ complex},\\
  -\operatorname{ord}_v(a)\log N(v),&v\text{ finite}.
 \end{cases}
\]
For a global function field with full constant field $\F_q$, set
$L_v(a)=-\deg(v)\operatorname{ord}_v(a)\log q$.
These conventions absorb the usual local weights into $L_v$; the product
formula is $\sum_vL_v(a)=0$.

\begin{theorem}[Rigidity of real global weights]\label{thm:global-weights}
Let $K$ be a global field and $(w_v)_v$ any family of real numbers indexed
by its standard places. Then
\[
 \sum_v w_vL_v(a)=0\quad\text{for every }a\in K^\times
 \quad\Longleftrightarrow\quad
 w_v=c\quad\text{for every }v
\]
for some $c\in\mathbb R$. Each tested sum is finite, so no boundedness
condition on the family is needed. In particular, a subset of the places
satisfies the restricted product formula exactly when it is empty or full.
\end{theorem}
\begin{proof}
Suppose first that $K$ is a number field. For a unit of its ring of
integers all finite logarithmic terms vanish. Dirichlet's unit theorem
says that the normalized archimedean logarithm vectors span the hyperplane
$\sum_{v\mid\infty}x_v=0$; see \citet[Theorem5.9]{MilneANT}.
The weight vector annihilates this hyperplane, hence every archimedean
weight has one value $c$.

Fix a finite prime ideal $\mathfrak p$. Finiteness of the ideal class group
\citep[Theorem4.4]{MilneANT} gives $m>0$ and $a\in K^\times$ with
$(a)=\mathfrak p^m$. All finite logarithmic terms of $a$ except the one at
$\mathfrak p$ vanish. Subtracting $c$ times the ordinary product formula
from the weighted identity gives
\[
 (w_{\mathfrak p}-c)L_{\mathfrak p}(a)=0.
\]
Since $L_{\mathfrak p}(a)=-m\log N(\mathfrak p)\ne0$, the remaining
weight is also $c$.

Now let $K$ be a global function field, and let $C/\F_q$ be its smooth
projective curve, of genus $g$. We first recall directly why its
degree-zero divisor class group is finite. Choose an effective divisor
$A$ of degree $N\ge g$. For every degree-zero divisor $D$, Riemann--Roch
\citep[Tag0B5B, Section53.5]{StacksRR} gives
\[
 \dim_{\F_q}H^0(C,\mathcal O_C(D+A))\ge N+1-g\ge1.
\]
Consequently $D$ is linearly equivalent to $E-A$ for an effective divisor
$E$ of degree $N$. Only finitely many such $E$ exist: their closed points
have degree at most $N$, and the points of each degree are finite because
$C(\F_{q^d})$ is finite. This proves the required class-group finiteness.

For any two places $v,w$, the divisor
$\deg(w)[v]-\deg(v)[w]$ has degree zero. A positive multiple is principal,
so choose $a\in K^\times$ and $m>0$ with
\[
 \operatorname{div}(a)
 =m\bigl(\deg(w)[v]-\deg(v)[w]\bigr).
\]
Its weighted logarithmic relation is
\[
 m\deg(v)\deg(w)\log q\,(w_w-w_v)=0.
\]
Thus $w_v=w_w$ for every pair. Conversely constant weights satisfy the
relation by the ordinary product formula. An indicator family is constant
exactly for the empty and full subsets, proving the final assertion.
\end{proof}

A weighted logarithmic height changes under scalar multiplication by the
weighted principal sum in Theorem~\ref{thm:global-weights}. Hence, on this
standard place system, it descends to projective space exactly for
constant weights. If one adjoins a separate measured family of trivial
absolute values, their logarithmic terms vanish identically and their
weights are unrestricted. The latter construction does not alter the
rigidity assertion for the original places.
